\documentclass[../../../include/open-logic-section]{subfiles}

\begin{document}

\olfileid{his}{set}{cantorplane}
\olsection{کانتور دربارهٔ خط و صفحه}

بخشی از اوضاع‌واحوال پیرامون اثبات شرودر--برنشتاین با تاریخی که در
\olref[his][set][pathology]{sec} بررسی کردیم پیوند دارد. به یاد آورید که
کانتور در ۱۸۷۷ ثابت کرد تعداد نقاط یک مربع دقیقاً با تعداد نقاط یکی از
ضلع‌های آن برابر است. در اینجا نخستین تلاش او برای اثبات را ارائه می‌کنیم.

بگذارید $\unitline$ خط واحد، یعنی مجموعهٔ نقاط $[0,1]$، و
$\unitsquare$ مربع واحد، یعنی مجموعهٔ نقاط
$\unitline \times \unitline$، باشد. با این اصطلاحات، کانتور ثابت کرد که
$\cardeq{\unitline}{\unitsquare}$. او یادداشتی برای ددکیند نوشت که در اصل
استدلال زیر را در برداشت.

\begin{thm}\ollabel{thm:cantorplane}
$\cardeq{\unitline}{\unitsquare}$
\end{thm}

\begin{proof}[اثبات: بخش نخست.]
$a,b\in\unitline$ را ثابت کنید. آن‌ها را در نمایش دودویی بنویسید؛ در این
صورت دنباله‌های نامتناهی‌ای از صفرها و یک‌ها، یعنی
$a_1,a_2,\dots$ و $b_1,b_2,\dots$، داریم به‌طوری‌که:
\begin{align*}
a &= 0.a_1a_2a_3a_4\dots\\
b &= 0.b_1b_2b_3b_4\dots
\intertext{اکنون تابع $f \colon \unitsquare \to \unitline$ را چنین در نظر بگیرید:}
f(a, b) & = 0.a_1b_1a_2b_2a_3b_3a_4b_4\dots
\end{align*}
تابع $f$ یک !!a{injection} است، زیرا اگر $f(a,b)=f(c,d)$، آنگاه برای هر
$n\in\Nat$ داریم $a_n=c_n$ و $b_n=d_n$؛ پس $a=c$ و $b=d$.
\end{proof}

متأسفانه، چنان‌که ددکیند به کانتور گوشزد کرد، این استدلال به پرسش اصلی
پاسخ نمی‌دهد. عدد
$0.\dot{1}\dot{0}=0.1010101010\ldots$ را در نظر بگیرید. باید داشته باشیم
$f(a,b)=0.\dot{1}\dot{0}$، که در آن:
\begin{align*}
a&=0.\dot{1}\dot{1}=0.111111\ldots\\
b&=0.
\end{align*}
اما $a=0.\dot{1}\dot{1}=1$ است. بنابراین وقتی می‌گوییم «$a$ و $b$ را در
نمایش دودویی بنویسید»، باید انتخاب کنیم از \emph{کدام} نمایش استفاده
می‌کنیم؛ و چون $f$ باید یک \emph{تابع} باشد، فقط می‌توانیم \emph{یکی} از
دو نمایش ممکن را به کار ببریم. اما اگر، برای مثال، نمایش متناهی را برگزینیم
و $a$ را به صورت «$1.000\ldots$» بنویسیم، هیچ زوج
$\tuple{a,b}$ای نداریم که $f(a,b)=0.\dot{1}\dot{0}$ باشد.

خلاصه اینکه ددکیند نشان داد با توجه به امکان بعضی بسط‌های دودوییِ متناوب،
تابع $f$ کانتور !!a{injection} است، اما \emph{!!a{surjection}} نیست. پس
کانتور تنها $\cardle{\unitsquare}{\unitline}$ را نشان داده است، نه
$\cardeq{\unitsquare}{\unitline}$ را.

کانتور تقریباً بی‌درنگ در پاسخ به ددکیند نوشت و در اصل پیشنهاد کرد که
اثبات را می‌توان چنین کامل کرد:

\begin{proof}[اثبات: تکمیل.]
پس نشان داده‌ایم $\cardle{\unitsquare}{\unitline}$. اما آشکارا یک
!!a{injection} از $\unitline$ به $\unitsquare$ وجود دارد: کافی است خط را
بر یکی از ضلع‌های مربع بخوابانیم. بنابراین
$\cardle{\unitline}{\unitsquare}$ و
$\cardle{\unitsquare}{\unitline}$. بنا بر قضیهٔ شرودر--برنشتاین
(\olref[sfr][siz][sb]{thm:schroder-bernstein})،
$\cardeq{\unitline}{\unitsquare}$.
\end{proof}

البته کانتور نمی‌توانست سطر آخر را با این اصطلاحات کامل کند، زیرا قضیهٔ
شرودر--برنشتاین هنوز اثبات نشده بود. گرچه کانتور بعداً آن را به صورت حدسی
کلی مطرح کرد، اثبات رضایت‌بخشی از آن تا ۱۸۹۷ به دست نیامد. (ازاین‌رو
کانتور بعدتر در همان سال ۱۸۷۷ اثبات دیگری برای
\olref{thm:cantorplane} عرضه کرد که از شرودر--برنشتاین استفاده نمی‌کرد.)

\end{document}
